%! Author = sriadityavedantam
%! Date = 3/30/26

\section{Field Extensions}

\begin{definition}[Vector Space]
    A vector space over $\ff$ is an additive abelian group $V$ equipped with scalar multiplication such that for $a,a_1,a_2\in\ff$ and $v,v_1,v_2\in V$, the following hold:
    \begin{enumerate}
        \item $a(v_1+v_2)=av_1+av_2$.
        \item $(a_1+a_2)v=a_{1}v+a_{2}v$.
        \item $a_1(a_{2}v)=(a_{1}a_2)v$.
        \item $1_\ff v=v$.
    \end{enumerate}
\end{definition}

\begin{definition}[Span]
    If every element of a vector space $V/\ff$ can be written as a linear combination of elements of a set $\{v_1,v_2,\hdots,v_n\}$, then we say that $\{v_1,v_2,\hdots,v_n\}$ spans $V/\ff$.
\end{definition}

\begin{definition}[Linearly Independent]
    A subset of a vector space $V/\ff$ is linearly independent over $\ff$ if whenever
    \[
        c_{1}v_1+\cdots+c_{n}v_n=0
    \]
    with $c_i\in\ff$, then $c_i=0_\ff$ for all $i$.
    Otherwise, it is dependent.
\end{definition}

\begin{definition}[Basis]
    A subset is called a basis if it is linearly independent and spans $V/\ff$.
\end{definition}

\begin{definition}[Dimension]
    If $p(x)\in\ff[x]$ is irreducible, then $\E$ is an extension field of $\ff$.
    In fact, this is a vector space over $\ff$.
    Its dimension is denoted by $[\E:\ff]$.
\end{definition}

\begin{theorem}{
    Suppose $K$ is an extension field of $\E$ and $\E$ is an extension field of $\ff$.
    Then
    \[
        [K:\ff]=[K:\E][\E:\ff].
    \]
}
Suppose
\[
    [\E:\ff]=n.
\]
Let $v_1,\hdots,v_n\in\E$ be a basis for $\E/\ff$.
Suppose
\[
    [K:\E]=m.
\]
Let $w_1,\hdots,w_m\in K$ be a basis for $K/\E$.
Our claim is that
\[
    \{w_{i}v_j:1\leq i\leq m,\ 1\leq j\leq n\}
\]
is a basis for $K/\ff$.
First, we show that $\{w_{i}v_j\}$ spans $K$.
Let $u\in K$.
Since $\{w_i\}$ spans $K/\E$, we can write
\[
    u=\sum_{i=1}^m \alpha_i w_i,
    \qquad \alpha_i\in\E.
\]
Since $\{v_j\}$ spans $\E/\ff$, each $\alpha_i$ can be written as
\[
    \alpha_i=\sum_{j=1}^n \beta_{ij}v_j,
    \qquad \beta_{ij}\in\ff.
\]
Therefore
\[
    u=\sum_{i=1}^m\sum_{j=1}^n \beta_{ij}w_{i}v_j.
\]
So $\{w_{i}v_j\}$ spans $K/\ff$.
Now suppose
\[
    \sum_{i=1}^m\sum_{j=1}^n \beta_{ij}w_{i}v_j=0,
    \qquad \beta_{ij}\in\ff.
\]
Rewrite this as
\[
    \sum_{i=1}^m \left(\sum_{j=1}^n \beta_{ij}v_j\right) w_i=0.
\]
Since the $w_i$ are linearly independent over $\E$, we have
\[
    \sum_{j=1}^n \beta_{ij}v_j=0
\]
for each $i$.
Since the $v_j$ are linearly independent over $\ff$, it follows that
\[
    \beta_{ij}=0
\]
for all $i,j$.
Thus $\{w_{i}v_j\}$ is linearly independent over $\ff$.
Hence it is a basis of $K/\ff$.
Therefore
\[
    [K:\ff]=mn=[K:\E][\E:\ff].
\]
\end{theorem}

\begin{definition}[Algebraic and Transcendental Functions]
    Let $\ff=\q$, $\E=\rr$, and $u=\pi$.
    There is no polynomial $p(u)=0$ with $p(x)\in\Q[x]$.
    If there is no such polynomial, we say $u$ is transcendental over $\ff$.
    If there is such a polynomial, we say $u$ is algebraic over $\ff$.
\end{definition}

To understand two versions of field extensions, let us look at when $\Q[\alpha]=\Q(\alpha)$.
We can use the first isomorphism theorem.

\begin{lemma}[$\Q[\alpha]=\Q(\alpha)$]{
    For the evaluation map $\phi:\Q[x]\to\Q[\alpha]$, we have
    \begin{align*}
        \text{$\phi$ is injective}
        &\iff \ker\phi=\{0\}\\
        &\iff \not\exists f(x)\in\Q[x]\text{ such that }f(\alpha)=0\\
        &\iff \text{$\alpha$ is transcendental over $\Q$.}
    \end{align*}
}
    Using the first isomorphism theorem, let $\phi$ be the homomorphism defined by evaluation at $\alpha$.
    Then
    \[
        \Ima\phi=\{f(\alpha):f(x)\in\Q[x]\}=\Q[\alpha].
    \]
    So $\phi$ is surjective.
    Also,
    \[
        \ker\phi=\{f(x)\in\Q[x]:f(\alpha)=0\}.
    \]
    Thus $\phi$ is injective if and only if $\ker\phi=\{0\}$, which happens if and only if there is no nonzero polynomial in $\Q[x]$ with $\alpha$ as a root.
    That is exactly the statement that $\alpha$ is transcendental over $\q$.
    Therefore, if $\alpha$ is transcendental over $\q$, then
    \[
        \Q[\alpha]\cong\Q[x].
    \]
    So $\Q[\alpha]$ is a ring, not a field.
\end{lemma}

\begin{definition}[Minimal Polynomial]
Suppose $p(x)$ is a monic polynomial of smallest degree such that
\[
p(\alpha)=0.
\]
This is called the minimal polynomial of $\alpha/\q$.
\end{definition}

\begin{lemma}{
The minimal polynomial $p(x)$ is irreducible.
}
Suppose $p(x)$ is reducible.
Then
\[
p(x)=q(x)g(x)
\]
for some nonconstant $q(x),g(x)\in\Q[x]$.
Evaluating at $\alpha$ gives
\[
0=p(\alpha)=q(\alpha)g(\alpha).
\]
So either $q(\alpha)=0$ or $g(\alpha)=0$.
But then either $q(x)$ or $g(x)$ is a polynomial of smaller degree than $p(x)$ having $\alpha$ as a root.
This contradicts the minimality of $p(x)$.
Therefore $p(x)$ is irreducible.
\end{lemma}

\begin{proof}
Using the first isomorphism theorem, suppose $\alpha$ is algebraic.
Let
\[
\phi:\Q[x]\to\Q[\alpha]
\]
be evaluation at $\alpha$.
Then
\[
\ker\phi=(p(x)),
\]
where $p(x)$ is the irreducible minimal polynomial of $\alpha$.
Therefore
\[
\Q[x]/(p(x))\cong\Q[\alpha].
\]
Since $p(x)$ is irreducible, $\Q[x]/(p(x))$ is a field.
So $\Q[\alpha]$ is a field.
Hence
\[
\Q[\alpha]=\Q(\alpha).
\]
\end{proof}

We have previously learned that $\Q(r)$ is a vector space.

{\Huge{$\Q(r)=\Q[r]$}}\\

What is the dimension $[\Q[r]:\Q]$.
We have to find the basis.
So what is the basis of $\Q[r]/\q$.

\begin{lemma}{
A basis is
\[
1,r,r^2,\hdots,r^{n-1},
\]
where $n=\deg f(x)$ and $f(x)$ is the minimal polynomial of $r/\q$.
}
    Since
    \[
        \Q[r]=\{g(r):g(x)\in\Q[x]\},
    \]
    every element of $\Q[r]$ is obtained by evaluating a polynomial at $r$.
    Let $f(x)$ be the minimal polynomial of $r$, with $\deg f(x)=n$ and $f(r)=0$.
    Take any $g(x)\in\Q[x]$.
    By the division algorithm,
    \[
        g(x)=f(x)q(x)+s(x),
    \]
    where $\deg s(x)<\deg f(x)=n$.
    Plug in $r$:
    \[
        g(r)=f(r)q(r)+s(r)=s(r),
    \]
    since $f(r)=0$.
    Thus $g(r)$ is a linear combination of
    \[
        1,r,r^2,\hdots,r^{n-1}.
    \]
    So these elements span $\Q[r]$.
    To see they are linearly independent, suppose
    \[
        c_0+c_1r+\cdots+c_{n-1}r^{n-1}=0
    \]
    with $c_i\in\Q$.
    Then the polynomial
    \[
        c_0+c_1x+\cdots+c_{n-1}x^{n-1}
    \]
    has $r$ as a root, but its degree is less than $n$.
    This contradicts the minimality of $f(x)$ unless all $c_i=0$.
    Therefore
    \[
        1,r,r^2,\hdots,r^{n-1}
    \]
    is a basis.
\end{lemma}

\begin{proof}

\end{proof}

\begin{definition}[Adjoin]
We say that $\mathbb{F}(u)$ is a field made by adjoining $u$ to $\mathbb{F}$.
\end{definition}

\begin{theorem*}{
Let $\mathbb{K}/\mathbb{F}$ and let $u\in\mathbb{K}$ be an algebraic element over $\mathbb{F}$ with minimal polynomial $p(x)$ of degree $n$.
Then
\begin{enumerate}
\item $\mathbb{F}(u)\cong \mathbb{F}[x]/(p(x))$;
\item $\{1_{\mathbb{F}},u,u^2,\ldots,u^{n-1}\}$ is a basis of the vector space $\mathbb{F}(u)$ over $\mathbb{F}$;
\item $[\mathbb{F}(u):\mathbb{F}]=n$.
\end{enumerate}
}
\end{theorem*}

\begin{corollary}{
If $u$ and $v$ have the same minimal polynomial $p(x)$ in $\mathbb{F}[x]$, then
\[
\mathbb{F}(u)\cong\mathbb{F}(v).
\]
}
\end{corollary}

\begin{definition}[Algebraic Extension]
An extension field $\mathbb{K}$ of a field $\mathbb{F}$ is said to be an algebraic extension if every element of $\mathbb{K}$ is algebraic over $\mathbb{F}$.
\end{definition}

\begin{theorem}{
If $\mathbb{K}$ is a finite-dimensional extension field of $\mathbb{F}$, then $\mathbb{K}$ is an algebraic extension of $\mathbb{F}$.
}
Let $u\in\mathbb{K}$.
Since $\mathbb{K}$ is finite-dimensional over $\mathbb{F}$, the set
\[
\{1,u,u^2,\hdots,u^n\}
\]
must be linearly dependent over $\mathbb{F}$ for $n$ large enough.
So there exist $a_0,\hdots,a_n\in\mathbb{F}$, not all zero, such that
\[
a_0+a_1u+\cdots+a_nu^n=0.
\]
Thus $u$ satisfies a nonzero polynomial with coefficients in $\mathbb{F}$.
So $u$ is algebraic over $\mathbb{F}$.
Since $u$ was arbitrary, every element of $\mathbb{K}$ is algebraic over $\mathbb{F}$.
Hence $\mathbb{K}$ is an algebraic extension of $\mathbb{F}$.
\end{theorem}